\section*{第 \ref{ch:g11:binom} 章 --- 二项分布}

\begin{solution}{exo:g11:binom:1}
第 $6$ 行与第 $7$ 行：
\[
1,\ 6,\ 15,\ 20,\ 15,\ 6,\ 1
\qquad\text{和}\qquad
1,\ 7,\ 21,\ 35,\ 35,\ 21,\ 7,\ 1 .
\]
因此 $\binom62 = 15$，$\binom73 = 35$，$\binom74 = 35$（对称性
$\binom73 = \binom74$ 体现成功与失败的互换）。
\end{solution}

\begin{solution}{exo:g11:binom:2}
试验次数固定（$4$ 次掷骰子），每次掷骰子有两种结果（点数为六或不为六，
$p = \frac16$），各次掷骰子相互独立：$X \sim \mathcal B(4, \frac16)$。
\[
\P(X=0) = \left(\frac56\right)^4 = \frac{625}{1296} \approx 0.48,
\qquad
\P(X=1) = 4 \times \frac16\left(\frac56\right)^3 = \frac{500}{1296}
\approx 0.39,
\]
\[
\P(X \geq 2) = 1 - \frac{625 + 500}{1296} = \frac{171}{1296}
\approx 0.13 .
\]
\end{solution}

\begin{solution}{exo:g11:binom:3}
\emph{1.} 为二项分布 $\mathcal B(20, \frac12)$：$n$ 固定，$p$ 相同，
各次抛掷相互独立。

\emph{2.} 不是二项分布：抽牌是\emph{不放回}地进行，
因此抽到 A 的概率逐次改变，各次抽取
并不相互独立。

\emph{3.} 由所给的独立性假设，
为二项分布 $\mathcal B(7, 0.3)$。
\end{solution}

\begin{solution}{exo:g11:binom:4}
$\E(X) = 50 \times 0.2 = 10$;
$\V(X) = 50 \times 0.2 \times 0.8 = 8$;
$\sigma(X) = 2\sqrt2 \approx 2.83$.
\end{solution}

\begin{solution}{exo:g11:binom:5}
$X \sim \mathcal B(6, 0.7)$.
\[
\P(X = 4) = \binom64 (0.7)^4 (0.3)^2 = 15 \times 0.2401 \times 0.09
\approx 0.324 .
\]
\[
\P(X \geq 5) = \binom65 (0.7)^5(0.3) + (0.7)^6
= 6 \times 0.16807 \times 0.3 + 0.117649 \approx 0.420 .
\]
\end{solution}

\begin{solution}{exo:g11:binom:6}
$X \sim \mathcal B\left(8, \frac12\right)$；每条路径的概率均为
$\frac{1}{256}$，于是
\[
\P(X \geq 6) = \frac{\binom86 + \binom87 + \binom88}{256}
= \frac{28 + 8 + 1}{256} = \frac{37}{256} \approx 0.14 .
\]
凭猜测通过的概率约为七分之一。
\end{solution}

\begin{solution}{exo:g11:binom:7}
“每种至少一个”的对立事件是“五个全部相同”：
全为 A 或全为 B，各自概率为 $\left(\frac12\right)^5 =
\frac1{32}$，于是
\[
\P(\text{每种各一个}) = 1 - \frac{2}{32} = \frac{15}{16} .
\]
\end{solution}

\begin{solution}{exo:g11:binom:8}
$\P(X = 3) = p^3$ 与 $\P(X \geq 1) = 1 - (1-p)^3$。解方程
$p^3 = \frac{27}{64} = \left(\frac34\right)^3$，得 $p = \frac34$（立方函数
严格递增，见\cref{ch:g11:func}，因此解
唯一）。
\end{solution}

\begin{solution}{exo:g11:binom:9}
$\P(\text{至少出现一次正面}) = 1 - \left(\frac12\right)^n$，因此
条件为 $\left(\frac12\right)^n < 0.01$，即 $2^n > 100$。由于
$2^6 = 64$ 与 $2^7 = 128$，故抛掷 $n = 7$ 次及以上即可。
\end{solution}

\begin{solution}{exo:g11:binom:10}
第 $n$ 行之和统计了 $n$ 次试验的树状图中全部路径，并按
成功次数分类。但树状图每经一次试验，路径数便加倍
（每条路径分裂为 S 与 F），因此共有 $2^n$ 条路径。于是
$\sum_{k} \binom nk = 2^n$。另外用归纳法：第 $0$ 行之和为
$1 = 2^0$，由帕斯卡法则，第 $n$ 行的每一项恰好对应第 $n+1$ 行中的两项，
故各行之和加倍。
\end{solution}

\begin{solution}{exo:g11:binom:11}
\emph{1.} 若该断言成立，则 $X \sim \mathcal B(10, 0.6)$。将
前几项相加：
\begin{align*}
\P(X \leq 3)
&= (0.4)^{10} + 10(0.6)(0.4)^9 + 45(0.6)^2(0.4)^8
+ 120(0.6)^3(0.4)^7\\
&\approx 0.0001 + 0.0016 + 0.0106 + 0.0425 = 0.0548 .
\end{align*}

\emph{2.} 与观察结果同样或更为极端的结果（至多 $3$ 次
赞成）的概率约为 $5.5\%$ --- 刚好\emph{高于} $5\%$
的临界值。严格依据规则，观察结果（勉强）
与该断言相容，故不拒绝它。此例说明
临界值的选取对临界决策影响有多大：若采用
$6\%$ 的约定，结论便会翻转。
\end{solution}

\begin{solution}{pb:g11:binom:1}
\textbf{1.} 各行依次为：$1$；$1\,1$；$1\,2\,1$；$1\,3\,3\,1$；
$1\,4\,6\,4\,1$; $1\,5\,10\,10\,5\,1$;
$1\,6\,15\,20\,15\,6\,1$。对称性：选出要取的 $k$ 个
对象，与选出要留下的 $n - k$ 个
对象，是同一个动作。

\textbf{2.} $1 + 4 + 6 + 4 + 1 = 16 = 2^4$；
$1 + 5 + 10 + 10 + 5 + 1 = 32 = 2^5$。证明：
$\binom nk$ 计数的是 $n$ 元集合各大小的
子集，子集总数为 $2^n$（每个元素取或不取，
各自独立）。

\textbf{3.} 选出 $k$ 人组成委员会，候选人为 $n + 1$ 人，其中一人
是佐伊：不含佐伊的有 $\binom nk$ 种（从其他人中选出全部
$k$ 人）；含佐伊的有 $\binom{n}{k-1}$ 种（
选她的 $k - 1$ 名同事）。共计：
$\binom nk + \binom{n}{k-1}$.

\textbf{4.} 帕斯卡三角形第 $7$ 行：$1\,7\,21\,35\,\dots$：$35$。
公式：$\frac{7 \times 6 \times 5}{3 \times 2 \times 1} =
35$.

\textbf{5.} $1 + 3 + 6 + 10 = 20 = \binom63$。逐级展开：
$\binom63 = \binom52 + \binom53 = \binom52 + \binom42 +
\binom43 = \binom52 + \binom42 + \binom32 + \binom33$ ---
每次应用帕斯卡法则，都会从阶梯形的和中
剥去一项。

\textbf{6.} 弹跳次数 $n$ 固定；每次弹跳都是
独立的伯努利试验，且 $p = \frac12$ 相同；槽位
编号即成功次数（向右弹跳的次数）；
三个方框均已打勾（见\cref{met:g11:binom:recognize}）：
$\mathcal B\!\left(n, \frac12\right)$.

\textbf{7.} 概率依次为 $\frac{1}{16}, \frac{4}{16},
\frac{6}{16}, \frac{4}{16}, \frac{1}{16}$ 对应槽位
$0, \dots, 4$：中心槽位 $2$ 最为拥挤。

\textbf{8.} 期望个数 $= 1024 \times
\binom{10}{k}/1024 = \binom{10}{k}$：中间槽位
$\binom{10}{5} = 252$ 个球；槽位 $7$：$\binom{10}{7} = 120$；
两端各槽位：$1$ 个球。中央高耸，向两侧
对称地降至细若游丝般的两端：这便是钟形。

\textbf{9.} $\E = np = 5$；$V = np(1 - p) = 2.5$；
$\sigma \approx 1.58$。在 $2\sigma$ 以内，槽位 $2$ 至 $8$
共装有
\[
\frac{45 + 120 + 210 + 252 + 210 + 120 + 45}{1024}
= \frac{1002}{1024} \approx 98\,\%
\]
的球 --- 远优于切比雪夫的
通用 $75\,\%$
（见\cref{pb:g11:stat:1}）：钟形分布高度集中。

\textbf{10.} $\E = 6$，$V = 10 \times 0.6 \times 0.4 = 2.4$，
$\sigma \approx 1.55$：堆积依旧保持钟形，但
顶峰滑向槽位 $6$ --- 倾斜的木板，就是把有偏
硬币的效应直观地显现出来。

\textbf{11.} 槽位编号是许多微小、独立、大小相同的随机推动的\emph{和}，
而此类和总会自行组织成钟形：多数推动相互抵消，
极端情形则需要全体一致。
身高、测量误差以及无数自然量同样是
许多微小独立效应的和，这正是同一轮廓
处处出现的原因；证明这一点的定理，
便是中心极限定理。

\textbf{12.} 横扫：一队赢下全部 $4$ 场：
$2 \times \left(\frac12\right)^4 = \frac18$.

\textbf{13.} 七场比赛要求前六场后比分为 $3$--$3$：
$\binom63 \left(\frac12\right)^6 = \frac{20}{64} =
\frac{5}{16}$.

\textbf{14.} 以 $5$ 场结束：胜者须赢下 $3$ 场，
且均出自前 $4$ 场，并赢下第 5 场：
$2 \times \binom43 \left(\frac12\right)^5 = \frac14$。以
$6$ 场结束：$2 \times \binom53 \left(\frac12\right)^6 = \frac{5}{16}$。
结果在 $4, 5, 6, 7$ 上的分布为：
$\frac18, \frac14, \frac{5}{16}, \frac{5}{16}$（总和为 $1$）。
期望长度：$4 \cdot \frac18 + 5 \cdot \frac14 + 6 \cdot
\frac{5}{16} + 7 \cdot \frac{5}{16} = 5.8125$ 场。六场与
七场的系列赛最有可能出现 --- 戏剧性
已内嵌于赛制之中。

\textbf{15.} 以 $4$ 场获胜：$0.6^4 = 0.1296$；以 $5$ 场获胜：
$\binom43\,0.6^3 \times 0.4 \times 0.6 = 0.2074$；以 $6$ 场获胜：
$\binom53\,0.6^3 \times 0.4^2 \times 0.6 = 0.2074$；以 $7$ 场获胜：
$\binom63\,0.6^3 \times 0.4^3 \times 0.6 = 0.1659$。合计：
约为 $0.710$：每场获胜概率为 $60\,\%$ 的队伍赢得 $71\,\%$ 的
系列赛 --- 赛制放大了这一优势。

\textbf{16.} 单场决赛：$60\,\%$。三场两胜：
$p^2 + 2p^2 q = 0.36 + 0.288 = 0.648$。阶梯
$60\,\% \to 65\,\% \to 71\,\%$ 随赛程延长而延续：
赛程越长，更多的比赛会平抑运气（大数定律的缩影），
因此长决赛能够筛出真正的强者 --- 这恰恰是
各类联赛所要兜售的卖点。

\textbf{17.} 公平硬币：$\E = 50$，$\sigma = \sqrt{25} = 5$；
$z = \frac{62 - 50}{5} = 2.4$：超出了 $2\sigma$
的约定 --- 这枚硬币值得深入调查。

\textbf{18.} $\P(X = 0) = 0.98^{50} \approx 0.364$；
$\P(X = 1) = 50 \times 0.02 \times 0.98^{49} \approx 0.372$;
$\P(X = 2) = \binom{50}{2} 0.02^2 \times 0.98^{48} \approx
0.186$。于是 $\P(X \geq 3) \approx 1 - 0.922 = 0.078$：约
$7.8\,\%$ --- 高于 $5\,\%$ 的临界值，故（暂时）尚不足以
拒绝该断言；若第二批又出现次品，
局面便会改变。

\textbf{19.} $\P(\text{至少赢一次}) = 1 - 0.999^{1000}
\approx 1 - 0.368 = 0.632$：一千张彩票，每张中奖率为千分之一，
得到的并非必然，而是 $63\,\%$。反复出现的
$0.368$ 其实是 $\frac1e$ 的伪装 --- 常数 $e$ 正式
登场要等到十二年级。

\textbf{20.} 识别条件：$n$ 固定、独立性、
$p$ 恒定 --- 当且仅当三者齐备时，才是二项分布。表格：帕斯卡
三角形第 $n$ 行。形状：钟形，当
$p = \frac12$ 时对称，否则发生平移。定位：中心为 $np$，
离散程度为 $\sqrt{np(1-p)}$ --- 正是决策所用的标准分数。
延伸：堆积所逼近的光滑钟形曲线，以及
解释大型钉板为何从不说谎的大数定律。
\end{solution}
